\section{Introduction}
\label{sec:introduction}

Recent advances in robotic systems have improved autonomy in perception, manipulation, and planning, yet robots still fail in real-world environments due to perception errors, action execution failures, partial observability, and dynamic changes~\cite{kaipa2016addressing,kaelbling1998planning,kurniawati2016online}. These failures often arise because the robot must select actions from an incomplete or ambiguous estimate of the current world state. In tasks such as tomato harvesting and waste sorting, uncertainty about object states or categories may remain even after sensing, making human feedback a useful source of task-relevant information.

Existing ways of incorporating humans into robot autonomy remain limited. Continuous-supervision and human-on-the-loop systems rely on operators to monitor execution and detect failures~\cite{goodrich2008human,dragan2013legibility,scharre2015introduction,nahavandi2017trusted}, which can increase workload and fatigue~\cite{chen2014human,endsley2018automation,wong2017workload}. Failure-triggered systems reduce monitoring but ask only after errors occur, while fixed or uncertainty-threshold query systems can ask too often~\cite{knowno2023,liang2024introspective,mullen2025lbap}. These approaches leave two coupled questions unresolved: \textit{when} should the robot ask a human, and \textit{what} should it ask?

We trigger human feedback from the robot's belief update rather than treating it as continuous supervision or failure recovery. After each action, the robot compares the observation with possible transition outcomes encoded in its symbolic belief. If multiple task-relevant successor states remain plausible, the system estimates uncertainty using entropy and asks about the symbolic hypothesis expected to reduce the ambiguity. This follows the broader goal of selective assistance, where robots request human input only when it is useful~\cite{sadigh2017active,baraglia2017efficient,leusmann2023understanding}.

This paper proposes a knowledge-based belief-space planning framework for proactive human querying. The robot maintains a symbolic knowledge base~\cite{davis1993knowledge} with STRIPS-style state and action structure~\cite{fikes1971strips}, while online decision making is performed using a POMDP formulation and POMCP-based planning~\cite{kaelbling1998planning,silver2010monte}. Instead of committing uncertain perceptual outputs as deterministic facts, the robot maintains a particle belief over symbolic possible worlds. After each observation, the system updates the belief over reachable frontier states, computes entropy-based confidence, and either commits the most likely state to the knowledge base or asks the human a targeted question before replanning.

\placeholderfig{Overview of the proposed framework. The final paper should show the closed loop among sensing, symbolic belief update, POMCP planning, human query, and knowledge-base commitment.}{1.9in}{fig:overview}

The framework differs from continuous supervision, failure-triggered intervention, always-query policies, and action-level query systems: the user is not asked to monitor every decision or choose robot actions, but instead answers targeted questions about uncertain symbolic state facts.

We evaluate the framework in waste sorting and tomato harvesting, including real-robot tomato experiments. We compare against no-query, always-query, and random-query baselines, analyze sensitivity to the confidence threshold $\tau$, and conduct a user study to examine whether selective state-level querying reduces subjective workload.

The contributions and findings of this paper are:
\begin{itemize}[leftmargin=*]
    \item We present an integrated robot system that determines both \textit{when} it is uncertain and \textit{what} it should ask the human during task execution, by combining symbolic knowledge, belief-space planning, entropy-based uncertainty estimation, and hypothesis-level feedback.
    \item We show that the proposed system substantially reduces the number of human queries compared with an always-query system while maintaining comparable task success, demonstrating that selective querying can preserve task performance without continuous supervision.
    \item We show through a user study that reducing unnecessary queries also reduces human workload, supporting the central goal of making uncertainty-aware robot assistance easier for people to use.
    \item We validate the system across simulated and real robotic domains and compare it against no-query, always-query, and random-query baselines to isolate the effect of deciding when and what to ask.
\end{itemize}
